% ECONOMICS_PROBLEM: {"id": "J42-574", "title": "Magnitude and sources of labor-market monopsony", "jel_code": "J42", "source_id": "OP-0023"}
% !TeX program = xelatex
\documentclass[11pt,a4paper]{article}
\usepackage[margin=23mm]{geometry}
\usepackage{fontspec}
\setmainfont{Latin Modern Roman}
\usepackage{amsmath,amssymb,mathtools}
\usepackage{xcolor,enumitem,booktabs,longtable,array,xurl,hyperref}
\definecolor{muted}{HTML}{53636C}
\hypersetup{colorlinks=true,urlcolor=blue,linkcolor=blue}
\setlength{\parindent}{0pt}
\setlength{\parskip}{5pt}
\newcommand{\R}{\mathbb R}
\newcommand{\E}{\mathbb E}
\newcommand{\N}{\mathbb N}
\newcommand{\ind}{\mathbf 1}
\newcommand{\Var}{\operatorname{Var}}
\newcommand{\Cov}{\operatorname{Cov}}
\newcommand{\supp}{\operatorname{supp}}
\DeclareMathOperator*{\argmin}{arg\,min}
\DeclareMathOperator*{\argmax}{arg\,max}
\newcommand{\entrymeta}[3]{{\sffamily\small\color{muted}\textbf{\texttt{#1}}\quad|\quad #2\par #3\par}}
\newcommand{\Problemref}[1]{\texttt{#1}}
\newcommand{\JELref}[2]{\texttt{#2} (JEL \texttt{#1})}
\begin{document}
\section*{Magnitude and sources of labor-market monopsony}
\textbf{Problem ID:} \texttt{J42-574}\quad \textbf{JEL:} \texttt{J42}\par
% BEGIN ECONOMICS STATEMENT
\label{problem:OP-0023}
\label{atlas:prob:L3}
\noindent\textbf{Original atlas ID:} L3 (entry 23). \textbf{Classification:} Editorial identification problem with an explicitly derived benchmark identity.

\noindent\textbf{Original atlas question:} How much wage-setting power do firms possess and what creates it?

\subsection*{Definitions and assumptions}
Fix a labor market, an employer $j$, and a wage interval $I\subset(0,\infty)$. In a static benchmark, let $n_j(w;\xi)$ be the differentiable labor supplied to $j$ when it posts wage $w$, holding competitors' offers, amenities, and the environment $\xi$ fixed. Require $n_j>0$ and $\partial n_j/\partial w>0$ on $I$, and define the residual supply elasticity $\varepsilon_j(w;\xi)=w(\partial n_j/\partial w)/n_j$. The firm maximizes $R_j(n)-wn$, where $R_j$ is differentiable net revenue excluding wage payments; define its marginal revenue product $m_j=R'_j(n_j)$. Assume an interior optimum and that raising $w$ affects profits only through the wage bill and labor supplied.
\subsection*{Formal research question}
Under precisely this benchmark, the first-order condition implies
\[
 m_j=w_j\left(1+\varepsilon_j^{-1}\right),
 \qquad 1-\frac{w_j}{m_j}=\frac{1}{1+\varepsilon_j}.
\]
Using applications, quits, recruitment, wage changes, and production information, identify or bound the joint distribution of $(\varepsilon_j,m_j,w_j)$ across firms and worker types. Then determine how this distribution changes under specified interventions on search costs, transport access, employer entry, bargaining, or information. State the observed law $P$ and the restrictions defining compatible models $\mathcal M(P)$, rather than inferring a unique elasticity from concentration.
\subsection*{Interpretation and scope}
The equations are benchmark identities under their assumptions, not universal formulas for any wage gap. Recruitment expenses, dynamic retention, hours choices, wage discrimination, bargaining, and product-market effects change the relevant optimality conditions. A firm-specific wage intervention must hold competing offers and amenities fixed or model their response. A quit elasticity alone need not equal the total labor-supply elasticity because recruitment also changes employment. Concentration measures depend on geographic and occupational market definitions and can correlate with productivity or amenities. To compare causes, introduce explicit policy counterfactuals and report interaction-sensitive contributions rather than a residual assigned to monopsony. Welfare evaluation must include worker amenities and unemployment as well as wages.
\subsection*{Source audit and status}
The identities follow directly from the stated optimization problem. [S1] verifies Manning's 2003 book through author-affiliated records; [S2] studies concentration; [S3] reviews firm-related inequality and models. These support the research setting without establishing that concentration equals wage-setting power. The identification and mechanism target is an editorial extension, not a named unresolved theorem. General 2026 open status is not certified by these sources.

\subsection*{References}
\begin{enumerate}[label=\textnormal{[S\arabic*]},leftmargin=*]
\item Alan Manning (2003). \emph{Monopsony in Motion: Imperfect Competition in Labor Markets}. Princeton University Press. \url{https://cep.lse.ac.uk/_new/People/person.asp?id=738}. \textit{Locator:} Primary publisher, working-paper, or author record: bibliographic information and abstract; full-text theorem verification is not claimed. \textit{Role:} Book-level framework for imperfect labor-market competition. Book identity verified in author-affiliated LSE records (also turn98search38); publisher access was blocked, so the author record is supplied.
\item José Azar, Ioana Marinescu, and Marshall Steinbaum (2022). \emph{Labor Market Concentration}. Journal of Human Resources 57(S), S167--S199. \url{https://www.nber.org/papers/w24147}. \textit{Locator:} Primary publisher, working-paper, or author record: bibliographic information and abstract; full-text theorem verification is not claimed. \textit{Role:} Concentration evidence; supports a proxy distinction rather than equivalence with monopsony.
\item David Card, Ana Rute Cardoso, Joerg Heining, and Patrick Kline (2018). \emph{Firms and Labor Market Inequality: Evidence and Some Theory}. Journal of Labor Economics 36(S1). \url{https://doi.org/10.1086/694153}. \textit{Locator:} Primary publisher, working-paper, or author record: bibliographic information and abstract; full-text theorem verification is not claimed. \textit{Role:} Firm heterogeneity, worker sorting, and imperfect-competition framework.
\end{enumerate}

\subsection*{Common conventions: Source mathematical conventions}

Unless an entry states otherwise, agent and firm sets are finite. State and action spaces are measurable, strategies and outcome maps are measurable, probability laws are specified, and expectations exist for the targets under discussion. Infinite-horizon discount factors lie in $(0,1)$ and discounted objectives are finite. Norms, tolerances and complexity input sizes must be specified for computational comparisons. Constants or primitives that appear locally are inputs, not empirical facts asserted by this catalogue.

\textbf{Identification.} A statistical model $m\in\mathcal M$ implies an observable law $P_m$ and target $\theta(m)$. At law $P$, its identified set is
\[
\Theta(P;\mathcal M)=\{\theta(m):m\in\mathcal M,\ P_m=P\}.
\]
Point identification holds when this set is a singleton, or equivalently when $P_{m_1}=P_{m_2}$ implies $\theta(m_1)=\theta(m_2)$ for all compatible models. Sharp bounds mean bounds attained, or approached when the boundary is not attained, by compatible models. Writing this set is a definition, not a solution: the substantive task is to establish informative restrictions and derive or estimate the set. An unrestricted-confounding model may give only trivial bounds.

\textbf{Inference.} Identification, estimability and valid uncertainty quantification are separate questions. Fix $\alpha\in(0,1)$. A confidence procedure $C_n$ over a declared law class $\mathcal P$ has asymptotic uniform coverage at level $1-\alpha$ when
\[
\liminf_{n\to\infty}\inf_{P\in\mathcal P}\Pr_P\{\theta(P)\in C_n\}\geq1-\alpha.
\]
For set-identified targets the coverage event must be defined separately, for example coverage of the full identified set. If an entry uses identification-indexed models instead of uniquely indexed laws, the same coverage convention is applied over those models. Pointwise limits do not establish uniform validity, and asymptotic validity does not guarantee finite-sample precision. Sample sizes, dependence structures and weak-identification sequences are part of each application.

\textbf{Counterfactuals.} $Y_i(a)$ is a potential outcome under a specified intervention, including its timing, implementation and versions. It is not $Y_i$ conditional on voluntarily choosing $a$. A full intervention vector permits interference; shorthand scalar treatment notation does not silently impose no interference. Assignment, selection, attrition, anticipation and measurement must be specified. A claim about an instrument's local effect is not automatically a population policy effect.

\textbf{Economic equilibrium.} A counterfactual model includes best responses, constraints, feasibility, market clearing, state transitions and expectations consistent with the stated information. When several equilibria exist, either a declared selector is an input or the target is set-valued. Comparisons across policy regimes must use the stated selection convention. Reduced-form relationships are not presumed invariant after an equilibrium-changing intervention.

\textbf{Optimization and welfare.} Optimization is over the stated feasible set. If attainment is not established, $\sup$ or $\inf$ and $\varepsilon$-optimal choices replace an asserted maximizer or minimizer. For example, with value $V=\sup_{a\in\mathcal A}W(a)<\infty$, an $\varepsilon$-optimal set is $\{a\in\mathcal A:W(a)\geq V-\varepsilon\}$ for $\varepsilon>0$. Benefits, costs, risk and health or environmental outcomes can be aggregated only using declared units and conversion weights. Interpersonal utility weights, the treatment of fiscal transfers and foreign profits, discounting, and the behavioral welfare criterion are explicit inputs. A comparative-static derivative additionally requires existence and appropriate differentiability of the selected counterfactual.

\textbf{Sources.} Local labels [S1], [S2], and so forth restart in every question. Locators identify the relevant abstract, model, section or result at the level actually checked; a bibliographic or abstract-level verification is not represented as inspection of an entire proof. Working papers are distinguished from later publications and changing versions. Source summaries are paraphrased. All URL checks and editorial formulations refer to the audit date stated above.
% END ECONOMICS STATEMENT
\end{document}
